\section{Loss}

model တွေကဘယ်လိုသင်ကြသလဲ? လူတွေလိုပဲ အမှားတွေလုပ်ရင်း သင်ယူကြတယ်။
လုပ်ကြည့်လိုက်၊ မှားလိုက်၊ မှန်အောင်ကြိုးစားလိုက်နဲ့ပေါ့။ trials and errors ပေါ့။
အဲဒီတော့ model တွေကိုလဲ မှားရင် မှားကြောင်းပြောပြရမယ်။ model ရဲ့တာဝန်က အမှားအနည်းဆုံး
ဖြစ်အောင် လုပ်ရမှာ။ ဒီမှာက model မှားရင် မှားတယ်ပြောရုံနဲ့ မလုံလောက်ဘူး။
ဘယ််လောက်မှားနေတယ်ဆိုတဲ့ ပမာဏကို ကိန်းဂဏန်းနဲ့ ပြောနိုင်မှ ပိုပြီးထိရောက်တာ။

မှားတဲ့ ပမာဏကို loss လို့ခေါ်တယ်။ အမှားတွေကို နည်းအောင်လုပ်တာက သင်္ချာနည်းနဲ့ဆိုရင်
loss တန်ဖိုးကိန်းဂဏန်းကို ပိုပြီးသေးသထက် သေးသွားအောင် လုပ်ပေးနေတာပဲ ဖြစ်တယ်။